% Solutions to Chapter 5, from lectures/L05/appendix/solutions: the README's prose, and every
% source file typeset straight from the repository.

\section{Chapter 5: Templates}
\label{sol:sec:c5}
Reference implementations for the exercises in \secref{c5:sec:exercises}. Each set is one program
in \file{lectures/L05/appendix/solutions}, built and run with \code{make} in its directory.

\note Try to solve the exercises yourself before looking at the solutions.

\subsection{Function template solution}
\label{sol:set:c5.1}
Exercises~\ref{c5:ex:1.1} and~\ref{c5:ex:1.2} (Exercise Set~1), in
\file{function_template/main.cpp}: the bit utility function templates \code{clear()} and
\code{toggle()}.

The function template solution implements two generic bit operations:
\begin{itemize}
  \item \code{clear(T& reg, std::uint8_t bit)}
  \item \code{toggle(T& reg, const Bits... bits)}
\end{itemize}
Both templates use \code{static_assert()} to verify that the register type is integral.

The solution demonstrates how parameter packs and compile-time type constraints can make bit
operations reusable for different integral register types.

\cppfile{lectures/L05/appendix/solutions/function_template/main.cpp}

\noindent Exercise Set~2 has no program to write, and so no reference implementation: the answers
to Exercise~\ref{c5:ex:2.1} are the subject of \secref{c5:sec:staticassert} and
\secref{c5:sec:whycompiletime}.

\subsection{Class template solution}
\label{sol:set:c5.3}
Exercises~\ref{c5:ex:3.1} and~\ref{c5:ex:3.2} (Exercise Set~3), in \file{class_template}: the class
template \code{driver::timer::Timer} in \file{include/driver/timer/timer.hpp}, with a primary stub
implementation and a specialization for \code{driver::timer::Type::Stm32}, and a demonstration
program in \file{source/main.cpp}.

The class template solution demonstrates compile-time selection of timer implementations using
template specialization.

\subsubsection{What is implemented}
\begin{itemize}
  \item \code{driver::timer::Type} enum class with values \code{Stm32} and \code{Stub}.
  \item Primary template \code{template<Type T = Type::Stub> class Timer final} for the stub timer
    implementation.
  \item Full specialization \code{template<> class Timer<Type::Stm32> final} for the STM32 timer
    implementation.
\end{itemize}

\subsubsection{Shared public interface}
Both timer variants provide the same public API:
\begin{itemize}
  \item \code{timeout_ms()}
  \item \code{isRunning()}
  \item \code{isInitialized()}
  \item \code{start()}
  \item \code{stop()}
  \item \code{toggle()}
  \item \code{tick()}
  \item \code{hasTimedOut()}
\end{itemize}

\subsubsection{Main demo program}
The solution demo uses:
\begin{itemize}
  \item A stub timer with a \code{500 ms} timeout.
  \item An STM32 timer with a \code{1500 ms} timeout.
\end{itemize}
Both timers are run for \code{2000 ms}. A message is printed every time a timer times out.

This shows how the same API can work for both stub and STM32 timer implementations without
runtime polymorphism (i.e.\ interfaces).

\cppfile{lectures/L05/appendix/solutions/class_template/include/driver/timer/timer.hpp}
\cppfile{lectures/L05/appendix/solutions/class_template/source/main.cpp}

\subsubsection{Answers to the reflection questions}
These answer Exercise~\ref{c5:ex:3.2}.

\task{a) Why is the timer type selected at compile time?}
\begin{itemize}
  \item The timer type is selected at compile time so the compiler generates concrete code for the
    chosen implementation.
  \item This avoids runtime overhead from virtual dispatch and allows the compiler to optimize each
    concrete timer class.
\end{itemize}

\task{b) What is a template specialization?}
\begin{itemize}
  \item A template specialization is a custom implementation of a template for a specific template
    argument.
  \item In this exercise, the primary \code{Timer} template provides the stub timer, and the
    specialization provides the STM32 timer.
\end{itemize}

\task{c) What advantages does this design provide compared to using inheritance and virtual
functions?}
\begin{itemize}
  \item It avoids runtime polymorphism overhead, such as vtable lookups and virtual function calls.
  \item Each implementation is compiled separately, allowing more aggressive optimization.
  \item The same public interface is preserved while behavior is selected at compile time.
\end{itemize}

\task{d) When would inheritance be a better choice?}
\begin{itemize}
  \item Inheritance is better when the concrete implementation must be selected at runtime.
  \item It is useful for plugin-style architectures or when many interchangeable implementations
    must be handled through a common base pointer/reference.
\end{itemize}

\task{e) What happens if the specialization for \code{Type::Stm32} is removed and a
\code{Timer<Type::Stm32>} instance is created?}
\begin{itemize}
  \item If the specialization is removed, the compiler will use the primary template
    implementation for \code{Timer<Type::Stm32>}.
  \item That means the instance would behave like the stub timer rather than a true STM32 timer.
\end{itemize}

\task{f) Why might a template specialization be preferred over a virtual interface in a small
microcontroller system?}
\begin{itemize}
  \item Because it removes runtime dispatch overhead and eliminates the need for a vtable.
  \item It keeps code paths explicit and allows the compiler to optimize the code for the target
    hardware.
  \item This is often better for resource-constrained microcontrollers where performance and code
    size matter.
\end{itemize}
